\documentclass[11pt,a4paper]{article}
\usepackage[margin=25mm]{geometry}
\usepackage{amsmath,amssymb,mathtools}
\usepackage{booktabs}
\usepackage{hyperref}

\newcommand{\dd}{\mathop{}\!\mathrm{d}}
\newcommand{\bxi}{\boldsymbol{\xi}}
\newcommand{\bx}{\mathbf{x}}

\title{Tensor-Product \(N\times N\) Gauss--Legendre Pixel Integration}
\author{Renderer Convergence Conjecture --- Experiment 1}
\date{\today}

\begin{document}
\maketitle

\section{Purpose and convention}

This note specifies the \(N\times N\) Gauss--Legendre rule used by the
Experiment~1 custom image integrator and suitable for a Riley pixel-box
integration mode.  It approximates the \emph{area average} of a scalar
texture/intensity field \(I(x,y)\) over one axis-aligned pixel
\[
    P=[x_0,x_0+\Delta x]\times[y_0,y_0+\Delta y].
\]
The desired camera value is
\begin{equation}
    \bar I_P =
    \frac{1}{\Delta x\Delta y}
    \int_{y_0}^{y_0+\Delta y}\int_{x_0}^{x_0+\Delta x}
        I(x,y)\,\dd x\,\dd y.
    \label{eq:pixel-average}
\end{equation}
For a square camera pixel, set \(\Delta x=\Delta y=h\).  The rule has
\(N^2\) evaluations per output pixel.  Here \(N\) means the number of nodes
\emph{per direction}, not the total number of samples.

\section{One-dimensional Gauss--Legendre rule}

Let \((\xi_i,w_i)\), \(i=1,\ldots,N\), be the standard Gauss--Legendre nodes
and weights on \([-1,1]\):
\begin{equation}
    \int_{-1}^{1} f(\xi)\,\dd\xi
    \approx \sum_{i=1}^N w_i f(\xi_i).
    \label{eq:reference-rule}
\end{equation}
It is exact for polynomials through degree \(2N-1\).  The nodes are the roots
of the Legendre polynomial \(P_N\), and the weights may be computed as
\begin{equation}
    w_i =
    \frac{2}{(1-\xi_i^2)[P_N'(\xi_i)]^2}.
    \label{eq:reference-weights}
\end{equation}
In practice, use a trusted Gauss--Legendre routine (for example
\texttt{numpy.polynomial.legendre.leggauss(N)}) or a precomputed table.

\section{Map the rule to one pixel}

Map each reference node independently to the physical pixel:
\begin{align}
    x_i &= x_0+\frac{\Delta x}{2}(1+\xi_i),\\
    y_j &= y_0+\frac{\Delta y}{2}(1+\xi_j).
    \label{eq:node-map}
\end{align}
The Jacobian for the two-dimensional map is
\(\Delta x\Delta y/4\).  Therefore the \emph{integral} is approximated by
\begin{equation}
    \int_P I(x,y)\,\dd x\,\dd y
    \approx
    \frac{\Delta x\Delta y}{4}
    \sum_{j=1}^{N}\sum_{i=1}^{N}
        w_iw_j I(x_i,y_j).
    \label{eq:pixel-integral-rule}
\end{equation}
Dividing by the pixel area gives the form to use for a camera pixel average:
\begin{equation}
    \boxed{
    \bar I_P^{(N)}
    =
    \sum_{j=1}^{N}\sum_{i=1}^{N}
        \widehat w_i\widehat w_j I(x_i,y_j),
    \qquad
    \widehat w_i=\frac{w_i}{2}.
    }
    \label{eq:pixel-average-rule}
\end{equation}
The normalised weights sum to one:
\[
    \sum_{i=1}^{N}\widehat w_i=1,
    \qquad
    \sum_{j,i}\widehat w_i\widehat w_j=1.
\]
Consequently a constant texture remains exactly constant, a useful
implementation check.

\section{Implementation recipe}

\begin{enumerate}
    \item Generate \(\xi_i,w_i\) once for the selected order \(N\).
    \item Form offsets from the pixel lower-left corner:
    \[
        d_i=\frac{h}{2}(1+\xi_i),\qquad
        \widehat w_i=\frac{w_i}{2}.
    \]
    \item Form the tensor product.  For every pair \((i,j)\), use sample
    offset \((d_i,d_j)\) and weight
    \[
        W_{ji}=\widehat w_j\widehat w_i
              =\frac{w_jw_i}{4}.
    \]
    \item For each output pixel, evaluate the texture at its \(N^2\) sample
    positions and accumulate \(\sum_{j,i}W_{ji}I_{ji}\).
\end{enumerate}

Equivalent pseudocode is:
\begin{verbatim}
xi, w = gauss_legendre(N)       # nodes/weights on [-1, 1]
offset = 0.5 * (xi + 1) * h
weight = 0.5 * w

value = 0
for j in range(N):
    for i in range(N):
        sample_xy = pixel_lower_left + (offset[i], offset[j])
        value += weight[i] * weight[j] * shade(sample_xy)
\end{verbatim}

\noindent The Experiment~1 Python implementation uses exactly this
construction: \texttt{offsets = 0.5*(points+1)*pixel\_size},
\texttt{wts\_norm = weights*0.5}, and the outer product
\texttt{wts\_norm[:, None] * wts\_norm[None, :]}.  It therefore produces a
pixel \emph{average}; no further division by \(h^2\) is required.

\section{Riley integration semantics}

For an alternative Riley pixel-box integrator, the quadrature must be applied
to the same continuous camera-pixel footprint as the current SSAA rule.
For each Gauss point:
\begin{enumerate}
    \item construct the corresponding camera ray/sample position;
    \item perform the normal Riley geometry/inverse-map and visibility work
          at that point;
    \item evaluate the selected shader there; and
    \item add the result with tensor-product weight \(W_{ji}\).
\end{enumerate}
The final resolved pixel is the weighted sum, not the arithmetic mean.  An
arithmetic mean is correct only for equal-weight SSAA samples.  If a sample is
outside the image/mesh, it must follow the existing renderer's background and
coverage convention; its quadrature weight must not be silently renormalised
unless that convention explicitly defines a conditional average over covered
area.

\section{Low-order checks}

\begin{center}
\begin{tabular}{ccl}
\toprule
\(N\) & nodes \(\xi_i\) & weights \(w_i\) \\
\midrule
1 & \(0\) & \(2\) \\
2 & \(\pm1/\sqrt{3}\) & \(1,1\) \\
3 & \(0,\ \pm\sqrt{3/5}\) & \(8/9,\ 5/9,\ 5/9\) \\
\bottomrule
\end{tabular}
\end{center}

Thus \(1\times1\) Gauss is the pixel-centre sample.  A \(2\times2\) rule uses
points at \(h(1\pm1/\sqrt3)/2\) from the lower-left corner and assigns every
point weight \(1/4\).  Higher orders have unequal weights and are therefore
not interchangeable with regular-grid SSAA.

\section{Accuracy and cost}

The tensor-product rule is polynomial-exact up to degree \(2N-1\) in each
coordinate.  For the Experiment~1 eggbox, the integrand is trigonometric, so
the rule is not finite-order exact but converges rapidly as \(N\) grows.  The
cost is \(N^2\) shader/geometry evaluations per output pixel.  Compare it
against the analytic eggbox box average to choose the lowest adequate order.
The existing convergence parameter for Gauss integration is consequently
\(S=N^2\), matching the total point count reported for SSAA.

\section{Computing nodes and weights for arbitrary \(N\)}

For a production renderer, precompute the one-dimensional rule once per
requested order \(N\), cache it, and form the two-dimensional tensor product
from it.  Do \emph{not} derive the two-dimensional nodes independently.
Two standard arbitrary-order constructions are suitable.

\subsection{Golub--Welsch construction}

The most direct robust approach is the Golub--Welsch algorithm.  Form the
symmetric \(N\times N\) Jacobi matrix
\begin{equation}
    J =
    \begin{bmatrix}
    0 & \beta_1 &  &  \\
    \beta_1 & 0 & \ddots & \\
     & \ddots & \ddots & \beta_{N-1}\\
     &  & \beta_{N-1} & 0
    \end{bmatrix},
    \qquad
    \beta_k=\frac{k}{\sqrt{4k^2-1}},
    \quad k=1,\ldots,N-1.
    \label{eq:golub-welsch-jacobi}
\end{equation}
If \(\lambda_i\) is the \(i\)th eigenvalue and
\(\mathbf{v}_i\) its unit eigenvector, then
\begin{equation}
    \boxed{
        \xi_i=\lambda_i,
        \qquad
        w_i=2\,[\mathbf{v}_i]_1^2.
    }
    \label{eq:golub-welsch-rule}
\end{equation}
The matrix is real symmetric, so use a symmetric tridiagonal eigensolver.
This produces ordered nodes and positive weights with excellent numerical
stability.

\subsection{Newton refinement of Legendre roots}

If an eigensolver is unavailable, solve for the roots of \(P_N\) directly.
Evaluate the Legendre polynomials using
\begin{align}
    P_0(x)&=1, &
    P_1(x)&=x,\\
    P_{m+1}(x)&=
    \frac{(2m+1)xP_m(x)-mP_{m-1}(x)}{m+1},
    \qquad m=1,\ldots,N-1.
    \label{eq:legendre-recurrence}
\end{align}
At a trial point \(x\), compute the derivative without differentiating the
recurrence:
\begin{equation}
    P_N'(x)
    =
    \frac{N\left(P_{N-1}(x)-xP_N(x)\right)}{1-x^2}.
    \label{eq:legendre-derivative}
\end{equation}
For the \(k\)th positive root, a useful initial guess is
\begin{equation}
    x_k^{(0)}
    =
    \cos\left(
        \frac{\pi(k-\tfrac14)}{N+\tfrac12}
    \right),
    \qquad
    k=1,\ldots,\left\lceil\frac{N}{2}\right\rceil.
    \label{eq:root-initial-guess}
\end{equation}
Apply Newton iteration until the update is below a floating-point tolerance:
\begin{equation}
    x_k^{(r+1)}
    =
    x_k^{(r)}
    -
    \frac{P_N(x_k^{(r)})}{P_N'(x_k^{(r)})}.
    \label{eq:newton-legendre}
\end{equation}
After convergence, assign
\begin{equation}
    \boxed{
        \xi_k=x_k,
        \qquad
        w_k=
        \frac{2}
        {(1-x_k^2)[P_N'(x_k)]^2}.
    }
    \label{eq:arbitrary-order-weight}
\end{equation}
Use symmetry to fill the paired negative root and identical weight:
\[
    \xi_{N+1-k}=-\xi_k,
    \qquad
    w_{N+1-k}=w_k.
\]
For odd \(N\), include the central node \(\xi=0\) only once.

\subsection{Required validation checks}

For any generated rule, verify numerically that
\[
    -1<\xi_i<1,\qquad
    w_i>0,\qquad
    \sum_{i=1}^{N}w_i=2.
\]
For the pixel-average rule, the corresponding two-dimensional weights must
sum to one:
\[
    \sum_{j=1}^{N}\sum_{i=1}^{N}
    \frac{w_iw_j}{4}=1.
\]
These checks catch incorrect Jacobian factors, root ordering errors, and
accidental use of integral weights where average weights are required.

\end{document}
